% TOP_PROBLEM: {"id": 8700008, "title": "Conjecture 1.9", "category_id": 3, "category": {"id": 3, "name": "graph_theory", "display_name": "Graph Theory", "description": "Problems involving graphs, networks, and their properties.", "slug": "graph-theory", "order_index": 3, "created_at": "2026-07-31T15:26:25.671Z"}, "status": "open", "problem_number": "AMR-086-0008", "set_id": 13}
% FIRST_POSTED: 2026-10-01
% ENTRY_KIND: statement_only
% UnsolvedMath ID 8700008; source catalogue code AMR-086-0008
% !TeX program = lualatex
% Definition and sources only; this file is not an LLM solution attempt.
\documentclass[11pt,a4paper]{article}
\usepackage[margin=25mm]{geometry}
\usepackage{fontspec}
\setmainfont{lmroman10-regular.otf}[BoldFont=lmroman10-bold.otf,
 ItalicFont=lmroman10-italic.otf,BoldItalicFont=lmroman10-bolditalic.otf]
\usepackage{amsmath,amssymb,mathtools}
\usepackage{unicode-math}
\setmathfont{latinmodern-math.otf}
\AtBeginDocument{\let\setminus\smallsetminus}
\usepackage{xurl,hyperref}
\hypersetup{colorlinks=true,urlcolor=blue}
\setcounter{secnumdepth}{0}
\setlength{\parindent}{0pt}
\setlength{\parskip}{5pt}
\setlength{\emergencystretch}{3em}
\begin{document}
\section*{Conjecture 1.9}\label{8700008}
{\small UnsolvedMath ID 8700008 · AMR-086-0008 · Graph Theory · AMR Open Problem Lists}
\subsection{Definitions and mathematical statement}
A Markov triple is a triple of positive integers satisfying
\[
a^2+b^2+c^2=3abc.
\]
Call it normalized when $1\le a\le b\le c$. The ordering removes permutations of the same coordinates while allowing equal entries. A Markov number is a positive integer occurring as the largest coordinate of at least one normalized Markov triple.

For each positive integer $m$, define
\[
\mathcal P(m)=\{(a,b)\in\mathbb Z^2:1\le a\le b\le m,
\ a^2+b^2+m^2=3abm\}.
\]
Is $|\mathcal P(m)|\le1$ for every $m$? Equivalently, whenever $m$ is a Markov number, does its value determine the other two coordinates of the normalized triple uniquely? This is a statement about all positive integer solutions, without a restriction on the size, primality, or factorization of $m$.

The equation is symmetric, so keeping the largest coordinate fixed without an ordering convention would introduce irrelevant exchanges of $a$ and $b$. The proposed uniqueness is exactly uniqueness of the ordered normalized pair. It does not say that a number occurs in only one arbitrary position among all triples.

For a negative answer one must exhibit the same $m$ with two genuinely different pairs in $\mathcal P(m)$. For a positive answer, checking any finite range of $m$ is insufficient: the assertion ranges over the unbounded set of all integers for which such a pair exists.
\subsection{Short English statement}
Does the largest entry of a normalized positive Markov triple uniquely determine the other two entries?
\subsection{Sources}
\begin{itemize}
\item[S1] ULAM.AI and the UnsolvedMath Contributors, supplied problems.json, record 8700008 (AMR-086-0008), AMR Open Problem Lists. Catalogue identity and source transcription; CC BY 4.0.
\url{https://huggingface.co/datasets/ulamai/UnsolvedMath}.
\item[S2] Waldschmidt - Open Diophantine Problems (2004), Conjecture 1.9. Original source indexed by AMR-086-0008.
\url{https://arxiv.org/abs/math/0312440}.
\end{itemize}
\end{document}
